\subsubsection{Python Type System}
Python is a dynamically typed language in which variable types are determined
at runtime rather than at the point of declaration.
Because a variable’s type is defined by the value assigned during execution,
its type can change at any point while the program runs.
Moreover, Python adopts a pure object-oriented design in which “everything is
an object,” making the notion of object compatibility a central issue in type
analysis.

In dynamic languages such as Python, object compatibility is governed by duck
typing.
Under duck typing, compatibility is determined not by an object’s nominal type
(e.g., int, str, Parrot) but by whether it provides the required methods or
attributes (such as fly() or len()).

\begin{listing}[t]
\begin{minted}[linenos, numbers=left, frame=lines, framesep=2mm, xleftmargin=10pt, fontsize=\footnotesize]{python}
class Parrot:
  def fly(self):
    print("Parrot flying")

class Airplane:
  def fly(self):
    print("Airplane flying")

class Whale:
  def swim(self):
    print("Whale swimming")

def lift_off(entity):
  entity.fly()

parrot = Parrot()
airplane = Airplane()
whale = Whale()

# prints `Parrot flying`
lift_off(parrot)   

# prints `Airplane flying`
lift_off(airplane) 

# Error : `'Whale' object has no attribute 'fly'`
lift_off(whale)    
\end{minted}
\caption{Example Python program demonstrating object compatibility determined by duck typing.}
\label{lst:duckTyping}
\end{listing}

\autoref{lst:duckTyping} illustrates this principle. As shown in lines 20–21, instances of
Parrot and Airplane—despite having no inheritance or structural
relationship—can be passed to the same function lift\_off() without causing any
errors.
This demonstrates that compatibility is determined by the presence of the fly()
method rather than by the concrete classes Parrot or Airplane.
In contrast, when a Whale object, which does not define fly(), is passed to the
same function (line 22), a TypeError is raised due to the missing attribute.

Thus, type determination in Python cannot rely solely on nominal types; it must
also incorporate structural types defined by the methods and attributes
available on an object.
However, because these methods and attributes can be dynamically added,
removed, or modified during execution, performing precise type analysis
statically—without running the program—becomes inherently challenging.

